\chapter{2006年北京卷（文）}
\section{单选题}
\begin{problem}
设集合 \(A = \{x \mid 2x + 1 < 3\}\)，\(B = \{x \mid -3 < x < 2\}\)，则 \(A \cap B\) 等于（\quad）

\choices
    {\(\{x\mid -3 < x < 1\}\)}
    {\(\{x\mid 1 < x < 2\}\)}
    {\( \{x \mid x > -3\} \)}
    {\( \{x \mid x < 1\} \)}
\end{problem}

\begin{answer}
A
\end{answer}

\begin{solution}
由 \(2x+1<3\)，得 \(x<1\)，所以 \(A=(-\infty,1)\). 又 \(B=(-3,2)\)，因此 \(A\cap B=(-3,1)\)，故选 A.
\end{solution}

\begin{problem}
函数 \( y = 1 + \cos x \) 的图象（\quad）

\choices
    {关于 \(x\) 轴对称}
    {关于 \(y\) 轴对称}
    {关于原点对称}
    {关于直线 \(x = \frac{\pi}{2}\) 对称}
\end{problem}

\begin{answer}
B
\end{answer}

\begin{solution}
因为 \(\cos(-x)=\cos x\)，所以函数 \(y=1+\cos x\) 是偶函数，其图象关于 \(y\) 轴对称，故选 B.
\end{solution}

\begin{problem}
若 \(\bs{a}\) 与 \(\bs{b} - \bs{c}\) 都是非零向量，则“\(\bs{a} \cdot \bs{b} = \bs{a} \cdot \bs{c}\)”是“\(\bs{a} \perp (\bs{b} - \bs{c})\)”的（\quad）

\choices
    {充分而不必要条件}
    {必要而不充分条件}
    {充分必要条件}
    {既不充分也不必要条件}
\end{problem}

\begin{answer}
C
\end{answer}

\begin{solution}
由题意，\(\bs{a}\cdot\bs{b}=\bs{a}\cdot\bs{c}\) 等价于 \(\bs{a}\cdot(\bs{b}-\bs{c})=0\). 由于 \(\bs{a}\) 与 \(\bs{b}-\bs{c}\) 都是非零向量，这又等价于 \(\bs{a}\perp(\bs{b}-\bs{c})\)，所以两者互为充分必要条件，故选 C.
\end{solution}

\begin{problem}
在 \(1\)，\(2\)，\(3\)，\(4\)，\(5\) 这五个数字组成的没有重复数字的三位数中，各位数字之和为奇数的共有（\quad）

\choices
    {\(36\)个}
    {\(24\)个}
    {\(18\)个}
    {\(6\)个}
\end{problem}

\begin{answer}
B
\end{answer}

\begin{solution}
五个数字中有三个奇数 \(1\)，\(3\)，\(5\) 和两个偶数 \(2,4\). 三位数的各位数字之和为奇数时，只可能有一个奇数和两个偶数，或三个奇数.

有一个奇数时，先选一个奇数和两个偶数，再排列，所得三位数有 \(\mathrm C_3^1\mathrm C_2^2\cdot3!=18\) 个；有三个奇数时，三个奇数全部选出并排列，有 \(3!=6\) 个. 因此共有 \(18+6=24\) 个，故选 B.
\end{solution}

\begin{problem}
已知 \( f(x)=\begin{cases}
(3-a)x-4a,&x<1,\\
\log_{a}x,&x\geqslant1,
\end{cases} \) 是 \( (-\infty,+\infty) \) 上的增函数，那么 \(a\) 的取值范围是（\quad）

\choices
    {\((1,+\infty)\)}
    {\((-\infty,3)\)}
    {\(\left[\frac{3}{5}, 3\right]\)}
    {\((1,3)\)}
\end{problem}

\begin{answer}
D
\end{answer}

\begin{solution}
要使对数有意义，必须有 \(a>0\) 且 \(a\ne1\). 当 \(x<1\) 时，函数为一次函数，要递增必须有 \(3-a>0\)，即 \(a<3\)；当 \(x\geqslant1\) 时，\(\log_a x\) 递增要求 \(a>1\).

在分界点 \(x=1\) 处，右侧函数值为 \(\log_a1=0\). 由于左段斜率为正，左侧函数值的上确界为 \(3-5a\)，故还需 \(3-5a\leqslant0\)，即 \(a\geqslant\frac35\). 与 \(a>1\) 及 \(a<3\) 合并，得 \(1<a<3\)，故选 D.
\end{solution}

\begin{problem}
如果 \(-1\)，\(a\)，\(b\)，\(c\)，\(-9\) 成等比数列，那么（\quad）

\choices
    {\(b=3,ac=9\)}
    {\(b=-3\)，\(ac=9\)}
    {\(b=3\)，\(ac=-9\)}
    {\(b=-3\)，\(ac=-9\)}
\end{problem}

\begin{answer}
B
\end{answer}

\begin{solution}
设等比数列的公比为 \(q\). 由首项和末项得 \(-q^4=-9\)，所以 \(q^4=9\). 于是中间项 \(b=-q^2=-3\)，并且 \(ac=(-q)(-q^3)=q^4=9\)，故选 B.
\end{solution}

\begin{problem}
设 \(A\)，\(B\)，\(C\)，\(D\) 是空间四个不同的点，在下列命题中，不正确的是（\quad）

\choices
    {若 \(AC\) 与 \(BD\) 共面，则 \(AD\) 与 \(BC\) 共面}
    {若 \(AC\) 与 \(BD\) 是异面直线，则 \(AD\) 与 \(BC\) 是异面直线}
    {若 \(AB = AC\)，\(DB = DC\) 则 \(AD = BC\)}
    {若 \(AB = AC\)，\(DB = DC\) 则 \(AD \perp BC\)}
\end{problem}

\begin{answer}
C
\end{answer}

\begin{solution}
若直线 \(AC\) 与 \(BD\) 共面，则四个点 \(A\)，\(B\)，\(C\)，\(D\) 都在同一平面内，从而直线 \(AD\) 与 \(BC\) 共面. 因此命题 A 正确.

若 \(AD\) 与 \(BC\) 共面，则存在一个平面同时包含直线 \(AD\) 和 \(BC\)，因而四个点 \(A\)，\(B\)，\(C\)，\(D\) 共面，进而直线 \(AC\) 与 \(BD\) 共面. 所以当 \(AC\) 与 \(BD\) 是异面直线时，\(AD\) 与 \(BC\) 不可能共面，只能是异面直线，命题 B 正确.

若 \(AB=AC\)，则点 \(A\) 在 \(BC\) 的垂直平分面上；若 \(DB=DC\)，则点 \(D\) 也在这个垂直平分面上. 因此直线 \(AD\) 垂直于 \(BC\)，命题 D 正确.

命题 C 不一定成立. 例如取 \(B=(-1,0,0)\)，\(C=(1,0,0)\)，\(A=(0,1,0)\)，\(D=(0,2,0)\)，则 \(AB=AC=\sqrt{2}\)，\(DB=DC=\sqrt{5}\)，但 \(AD=1\ne2=BC\). 所以不正确的是 C.
\end{solution}

\begin{problem}
如图为某三岔路口交通环岛的简化模型，在某高峰时段，单位时间进出路口 \(A\)，\(B\)，\(C\) 的机动车辆数如图所示. 图中 \(x_{1}\)，\(x_{2}\)，\(x_{3}\) 分别表示该时段单位时间通过路段 \(\overline{AB}\)，\(\overline{BC}\)，\(\overline{CA}\) 的机动车辆数（假设：单位时间内，在上述路段中，同一路段上驶入与驶出的车辆数相等），则 \(x_{1}\)，\(x_{2}\)，\(x_{3}\) 的大小关系是（\quad）

\bitmapfigure[max width=0.42\linewidth]{img/2006/beijing_liberal/q8_fig1.png}

\choices
    {\(x_{1} > x_{2} > x_{3}\)}
    {\(x_{1} > x_{3} > x_{2}\)}
    {\(x_{2} > x_{3} > x_{1}\)}
    {\(x_{3} > x_{2} > x_{1}\)}
\end{problem}

\begin{answer}
C
\end{answer}

\begin{solution}
按图中车辆行驶方向，对路口 \(A\) 列车辆守恒方程，得 \(50+x_{3}=55+x_{1}\)，所以 \(x_{3}=x_{1}+5\). 对路口 \(B\) 列方程，得 \(x_{1}+30=20+x_{2}\)，所以 \(x_{2}=x_{1}+10\). 因此 \(x_{2}>x_{3}>x_{1}\)，故选 C.
\end{solution}

\section{填空题}
\begin{problem}
若三点 \(A(2,2)\)，\(B(a,0)\)，\(C(0,4)\) 共线，则 \(a\) 的值等于\fillinblank{}.
\end{problem}

\begin{answer}
\(4\)
\end{answer}

\begin{solution}
直线 \(AC\) 的斜率为 \(\dfrac{2-4}{2-0}=-1\)，所以其方程为 \(y=4-x\). 点 \(B(a,0)\) 在此直线上，故 \(0=4-a\)，解得 \(a=4\).
\end{solution}

\begin{problem}
在 \(\left(x - \frac{2}{x}\right)^7\) 的展开式中，\(x^3\) 的系数是\fillinblank{}.（用数字作答）
\end{problem}

\begin{answer}
\(84\)
\end{answer}

\begin{solution}
展开式的第 \(k+1\) 项为 \(\mathrm C_7^k x^{7-k}\left(-\dfrac{2}{x}\right)^k=\mathrm C_7^k(-2)^k x^{7-2k}\). 要得到 \(x^3\)，需满足 \(7-2k=3\)，所以 \(k=2\). 因此 \(x^3\) 的系数为 \(\mathrm C_7^2(-2)^2=84\)，故填 \(84\).
\end{solution}

\begin{problem}
已知函数 \( f(x) = a^x - 4a + 3 \) 的反函数的图象经过点 \((-1,2)\)，那么 \( a \) 的值等于\fillinblank{}.
\end{problem}

\begin{answer}
\(2\)
\end{answer}

\begin{solution}
反函数的图象与原函数图象关于直线 \(y=x\) 对称，所以原函数的图象经过点 \((2,-1)\). 代入 \(f(x)=a^x-4a+3\)，得 \(a^2-4a+3=-1\)，即 \((a-2)^2=0\). 又指数函数要求 \(a>0\) 且 \(a\ne1\)，故 \(a=2\).
\end{solution}

\begin{problem}
已知向量 \(\bs{a} = (\cos \alpha, \sin \alpha)\)，\(\bs{b} = (\cos \beta, \sin \beta)\)，且 \(\bs{a} \neq \pm \bs{b}\)，那么 \(\bs{a} + \bs{b}\) 与 \(\bs{a} - \bs{b}\) 的夹角的大小是\fillinblank{}.
\end{problem}

\begin{answer}
\(\dfrac{\pi}{2}\)
\end{answer}

\begin{solution}
由 \(\bs{a}\) 与 \(\bs{b}\) 的坐标表示，得 \(|\bs{a}|=|\bs{b}|=1\). 因此 \((\bs{a}+\bs{b})\cdot(\bs{a}-\bs{b})=|\bs{a}|^2-|\bs{b}|^2=0\). 题设 \(\bs{a}\ne\pm\bs{b}\) 保证了 \(\bs{a}+\bs{b}\) 与 \(\bs{a}-\bs{b}\) 都是非零向量，所以它们互相垂直，夹角为 \(\dfrac{\pi}{2}\).
\end{solution}

\begin{problem}
在 \(\triangle ABC\) 中，\(\angle A\)，\(\angle B\)，\(\angle C\) 所对的边长分别为 \(a\)，\(b\)，\(c\) 若 \(\sin A:\sin B:\sin C = 5:7:8\)，则 \(a:b:c =\fillinblank{}\)，\(\angle B\) 的大小是\fillinblank{}.
\end{problem}

\begin{answer}
\(5:7:8\)，\(\dfrac{\pi}{3}\)
\end{answer}

\begin{solution}
由正弦定理 \(\dfrac{a}{\sin A}=\dfrac{b}{\sin B}=\dfrac{c}{\sin C}\)，得 \(a:b:c=\sin A:\sin B:\sin C=5:7:8\).

由余弦定理，\(\cos B=\dfrac{a^2+c^2-b^2}{2ac}=\dfrac{5^2+8^2-7^2}{2\cdot5\cdot8}=\dfrac12\). 因为 \(B\) 是三角形内角，\(0<B<\pi\)，所以 \(B=\dfrac{\pi}{3}\). 故填 \(5:7:8\)，\(\dfrac{\pi}{3}\).
\end{solution}

\begin{problem}
已知点 \(P(x,y)\) 的坐标满足条件 \(\begin{cases}
x+y \leqslant 4,\\
y \geqslant x,\\
x \geqslant 1,
\end{cases}\) 点 \(O\) 为坐标原点，那么 \(|PO|\) 的最小值等于\fillinblank{}，最大值等于\fillinblank{}.
\end{problem}

\begin{answer}
\(\sqrt{2}\)，\(\sqrt{10}\)
\end{answer}

\begin{solution}
由 \(x\geqslant1\)、\(y\geqslant x\) 和 \(x+y\leqslant4\)，可知可行域是以 \((1,1)\)、\((1,3)\)、\((2,2)\) 为顶点的三角形.

令 \(r^2=|PO|^2=x^2+y^2\). 在可行域内 \(x\geqslant1\)、\(y\geqslant x\)，所以 \(r^2\geqslant1^2+1^2=2\)，且点 \((1,1)\) 可取到，故 \(|PO|_{\min}=\sqrt{2}\). 可行域是闭三角形，\(r^2\) 在其上能取得最大值；若最大值在内部取得，则应有 \(2x=0\) 且 \(2y=0\)，这与 \(x\geqslant1\)、\(y\geqslant1\) 矛盾，因此最大值在边界上取得. 再分别考察三条边：在 \(x=1\) 上，\(1\leqslant y\leqslant3\)，有 \(r^2=1+y^2\)，最大值为 \(1+3^2=10\)；在 \(y=x\) 上，\(1\leqslant x\leqslant2\)，有 \(r^2=2x^2\)，最大值为 \(2\cdot2^2=8\)；在 \(x+y=4\) 上，\(1\leqslant x\leqslant2\)，有 \(r^2=x^2+(4-x)^2=2(x-2)^2+8\)，最大值为 \(10\). 故可行域内 \(r^2\) 的最大值为 \(10\)，从而 \(|PO|_{\max}=\sqrt{10}\).
\end{solution}

\section{解答题}
\begin{problem}
已知函数 \( f(x) = \frac{1 - \sin 2x}{\cos x} \).

\begin{enumerate}
\item 求 \( f(x) \) 的定义域；
\item 设 \(\alpha\) 是第四象限的角，且 \(\tan \alpha = -\frac{4}{3}\)，求 \(f(\alpha)\) 的值.
\end{enumerate}
\end{problem}

\begin{solution}
\begin{enumerate}
\item 分式有意义的条件是 \(\cos x\ne0\). 因此 \(x\ne k\pi+\dfrac{\pi}{2}\)，其中 \(k\in\Z\)，故定义域为 \(\{x\mid x\ne k\pi+\dfrac{\pi}{2},\ k\in\Z\}\).
\item 因为 \(\alpha\) 是第四象限的角且 \(\tan\alpha=-\dfrac43\)，所以 \(\sin\alpha=-\dfrac45\)，\(\cos\alpha=\dfrac35\). 于是
\(f(\alpha)=\dfrac{1-\sin2\alpha}{\cos\alpha}=\dfrac{1-2\sin\alpha\cos\alpha}{\cos\alpha}=\dfrac{1-2\left(-\dfrac45\right)\left(\dfrac35\right)}{\dfrac35}=\dfrac{49}{15}\).
\end{enumerate}
\end{solution}

\begin{problem}
已知函数 \( f(x) = ax^3 + bx^2 + cx \) 在点 \( x_0 \) 处取得极大值 \(5\)，其导函数 \( y = f'(x) \) 的图象经过点 \((1,0)\)，\((2,0)\). 如图所示，求：

\begin{enumerate}
\item 求 \(x_0\) 的值；
\item 求 \(a\)，\(b\)，\(c\) 的值.
\end{enumerate}

\bitmapfigure[max width=0.28\linewidth]{img/2006/beijing_liberal/q16_fig1.png}
\end{problem}

\begin{solution}
\begin{enumerate}
\item 由图象可知，\(f'(x)\) 在 \((-\infty,1)\) 和 \((2,+\infty)\) 上为正，在 \((1,2)\) 上为负. 因此 \(f(x)\) 在 \(x=1\) 处由增变减，取得极大值，故 \(x_0=1\).
\item 因为 \(f'(x)=3ax^2+2bx+c\)，且图象与 \(x\) 轴交于 \(1,2\)，又在两根外侧为正，所以可设 \(f'(x)=m(x-1)(x-2)\)，其中 \(m>0\). 比较系数得 \(3a=m\)，\(2b=-3m\)，\(c=2m\)，即 \(b=-\dfrac92a\)，\(c=6a\). 又 \(f(1)=5\)，所以 \(a+b+c=5\)，即 \(a-\dfrac92a+6a=5\)，解得 \(a=2\)，从而 \(b=-9\)，\(c=12\).
\end{enumerate}
\end{solution}

\begin{problem}
如图，\(ABCD - A_{1}B_{1}C_{1}D_{1}\) 是正四棱柱.

\begin{enumerate}
\item 求证：\(BD \perp\) 平面 \(ACC_{1}A_{1}\)；
\item 若二面角 \(C_{1} - BD - C\) 的大小为 \(60^{\circ}\)，求异面直线 \(BC_{1}\) 与 \(AC\) 所成角的大小.
\end{enumerate}

\bitmapfigure[max width=0.45\linewidth]{img/2006/beijing_liberal/q17_fig1.png}
\end{problem}

\begin{solution}
\begin{enumerate}
\item 设正方形底面的边长为 \(s\)，棱柱的高为 \(h\). 在正方形 \(ABCD\) 中，\(BD\perp AC\)；又 \(BD\perp CC_1\)，且 \(AC\) 与 \(CC_1\) 是平面 \(ACC_1A_1\) 内相交的两条直线. 因此 \(BD\perp\) 平面 \(ACC_1A_1\).
\item 设 \(O=AC\cap BD\). 由上问，\(BD\perp AC\)，又 \(BD\perp CC_1\)，所以 \(BD\perp C_1O\). 因此二面角 \(C_1-BD-C\) 的平面角为 \(\angle C_1OC=60^{\circ}\). 在直角三角形 \(C_1OC\) 中，\(OC=\dfrac{s}{\sqrt2}\)，所以 \(\tan60^{\circ}=\dfrac{CC_1}{OC}=\dfrac{h}{s/\sqrt2}\)，从而 \(h^2=\dfrac32s^2\).

建立直角坐标系，取 \(A\) 为原点，\(AB\)，\(AD\)，\(AA_1\) 分别为三个坐标轴正向，则 \(B=(s,0,0)\)，\(C=(s,s,0)\)，\(C_1=(s,s,h)\). 因此直线 \(BC_1\) 的方向向量可取 \((0,s,h)\)，直线 \(AC\) 的方向向量可取 \((s,s,0)\). 设两直线所成的锐角为 \(\theta\)，则
\(\cos\theta=\dfrac{s^2}{\sqrt{s^2+h^2}\cdot s\sqrt2}=\dfrac{1}{\sqrt{2\left(1+h^2/s^2\right)}}=\dfrac{1}{\sqrt5}=\dfrac{\sqrt5}{5}\).
所以 \(\theta=\arccos\dfrac{\sqrt5}{5}\).
\end{enumerate}
\end{solution}

\begin{problem}
某公司招聘员工，指定三门考试课程，有两种考试方案. 方案一：考试三门课程，至少有两门及格为考试通过； 方案二：在三门课程中，随机选取两门，这两门都及格为考试通过. 假设某应聘者对三门指定课程考试及格的概率分别是 \(0.5\)，\(0.6\)，\(0.9\)，且三门课程考试是否及格相互之间没有影响.

\begin{enumerate}
\item 该应聘者用方案一考试通过的概率；
\item 该应聘者用方案二考试通过的概率.
\end{enumerate}
\end{problem}

\begin{solution}
\begin{enumerate}
\item 三门都及格的概率为 \(0.5\times0.6\times0.9=0.27\). 恰有两门及格包括第一、二门及格且第三门不及格，第一、三门及格且第二门不及格，以及第二、三门及格且第一门不及格，其概率为
\(0.5\times0.6\times\left(1-0.9\right)+0.5\times\left(1-0.6\right)\times0.9+\left(1-0.5\right)\times0.6\times0.9=0.03+0.18+0.27=0.48\). 所以方案一通过的概率为 \(0.27+0.48=0.75\).
\item 随机选取的两门可能是第一、二门，第二、三门，第一、三门，三种情况等可能. 因此方案二通过的概率为
\(\dfrac13(0.5\times0.6+0.6\times0.9+0.5\times0.9)=\dfrac{1.29}{3}=0.43\).
\end{enumerate}
\end{solution}

\begin{problem}
椭圆 \(C: \frac{x^2}{a^2} + \frac{y^2}{b^2} = 1\)（\(a > b > 0\)）的两个焦点为 \(F_1\)，\(F_2\)，点 \(P\) 在椭圆 \(C\) 上. 且 \(PF_1 \perp F_1F_2\)，\(|PF_1| = \frac{4}{3}\)，\(|PF_2| = \frac{14}{3}\).

\begin{enumerate}
\item 求椭圆 \(C\) 的方程；
\item 若直线 \(l\) 过圆 \(x^{2} + y^{2} + 4x - 2y = 0\) 的圆心 \(M\)，交椭圆 \(C\) 于 \(A\)，\(B\) 两点，且 \(A\)，\(B\) 关于点 \(M\) 对称，求直线 \(l\) 的方程.
\end{enumerate}
\end{problem}

\begin{solution}
\begin{enumerate}
\item 设椭圆的半焦距为 \(c\). 取焦点为 \(F_1(-c,0)\)、\(F_2(c,0)\). 由 \(PF_1\perp F_1F_2\)，点 \(P\) 的横坐标为 \(-c\)，且 \(|PF_1|=\dfrac43\). 又由椭圆的定义，
\(2a=|PF_1|+|PF_2|=\dfrac43+\dfrac{14}{3}=6\)，所以 \(a=3\). 在直角三角形 \(PF_1F_2\) 中，\((2c)^2+\left(\dfrac43\right)^2=\left(\dfrac{14}{3}\right)^2\)，解得 \(c^2=5\). 于是 \(b^2=a^2-c^2=4\)，椭圆方程为 \(\dfrac{x^2}{9}+\dfrac{y^2}{4}=1\).
\item 将圆方程化为 \((x+2)^2+(y-1)^2=5\)，所以圆心 \(M=(-2,1)\). 设直线 \(l\) 过 \(M\)，方向向量为 \((u,v)\)，其上的点可表示为 \((x,y)=(-2+tu,1+tv)\). 代入椭圆方程，关于 \(t\) 的二次方程的一次项系数为 \(-\dfrac{4u}{9}+\dfrac{v}{2}\).

因为交点 \(A\)，\(B\) 关于 \(M\) 对称，所以对应的两个参数根互为相反数，一次项系数必须为零，即 \(-\dfrac{4u}{9}+\dfrac{v}{2}=0\)，从而 \(v=\dfrac89u\). 因此直线斜率为 \(\dfrac89\)，故 \(y-1=\dfrac89(x+2)\)，即 \(8x-9y+25=0\).
\end{enumerate}
\end{solution}

\begin{problem}
设等差数列 \(\{a_{n}\}\) 的首项 \(a_1\) 及公差 \(d\) 都为整数. 前 \(n\) 项和为 \(S_{n}\).

\begin{enumerate}
\item 若 \(a_{11} = 0\)，\(S_{14} = 98\)，求数列 \(\{a_n\}\) 的通项公式；
\item 若 \(a_1 \geqslant 6\)，\(a_{11} > 0\)，\(S_{14} \geqslant 77\)，求所有可能的数列 \(\{a_n\}\) 的通项公式.
\end{enumerate}
\end{problem}

\begin{solution}
\begin{enumerate}
\item 由 \(a_{11}=a_1+10d=0\)，得 \(a_1=-10d\). 又
\(S_{14}=\dfrac{14}{2}(2a_1+13d)=98\)，所以 \(2a_1+13d=14\). 代入 \(a_1=-10d\)，得 \(-7d=14\)，即 \(d=-2\)，\(a_1=20\). 因此
\(a_n=a_1+(n-1)d=20-2(n-1)=22-2n\)，其中 \(n=1\)，\(2\)，\(3\)，\(\ldots\).
\item 由 \(a_1\)，\(d\) 为整数，条件 \(a_{11}>0\) 等价于 \(a_1+10d\geqslant1\)，而
\(S_{14}=7(2a_1+13d)\geqslant77\) 等价于 \(2a_1+13d\geqslant11\). 所以所有可能的整数对 \((a_1,d)\) 恰好满足 \(a_1\geqslant6\)，\(a_1+10d\geqslant1\)，\(2a_1+13d\geqslant11\).

等价地，对任意整数 \(d\)，取整数 \(a_1\geqslant\max\left\{6,1-10d,\left\lceil\dfrac{11-13d}{2}\right\rceil\right\}\)，就得到一个符合条件的数列. 按 \(d\) 的取值分类，这也就是：当 \(d\geqslant0\) 时 \(a_1\geqslant6\)；当 \(d=-1\) 时 \(a_1\geqslant12\)；当 \(d\leqslant-2\) 时 \(a_1\geqslant1-10d\). 所有可能的通项公式为
\(a_n=a_1+(n-1)d\)，其中 \(n=1\)，\(2\)，\(3\)，\(\ldots\).
\end{enumerate}
\end{solution}
